% Chapter 7 - Requirements
% Neues Kapitel (Rueckmeldung des Betreuers). Nimmt Problem Statement,
% Objectives und Scope aus der (jetzt kurzen) Einleitung auf; ergaenzt um
% die technischen Anforderungen (eigener Vorschlag, siehe Strukturplan.tex).
\section{Requirements}
\label{chap:requirements}

\subsection{Problem Statement}
\label{sec:problemStatement}
% TODO: siehe Strukturplan.tex, Abschnitt 7.1
% - McEliece ist mathematisch anspruchsvoll
% - Vorhandene Literatur richtet sich haeufig an Experten
% - Fehlende interaktive Lernwerkzeuge

\subsection{Objectives}
\label{sec:objectives}
% TODO: siehe Strukturplan.tex, Abschnitt 7.2
% - Entwicklung eines CTO-Plugins, Visualisierung zentraler Ablaeufe, Unterstuetzung des Lernprozesses

\subsection{Scope}
\label{sec:scope}
% TODO: siehe Strukturplan.tex, Abschnitt 7.3
% - Abgrenzung der Arbeit; ggf. Festlegung Lehr-/Demo-Parameter vs. reale
%   Sicherheitsparameter (siehe eigene Anmerkung in Strukturplan.tex)

\subsection{Functional Requirements}
\label{sec:functionalRequirements}
% TODO: siehe Strukturplan.tex, Abschnitt 7.4

\subsection{Non-Functional Requirements}
\label{sec:nonFunctionalRequirements}
% TODO: siehe Strukturplan.tex, Abschnitt 7.5

\subsection{Requirements for the McEliece Plugin}
\label{sec:pluginRequirements}
% TODO: siehe Strukturplan.tex, Abschnitt 7.6
% - Fachliche Anforderungen, Didaktische Anforderungen, Technische Anforderungen
